\documentclass[11pt]{article}
\usepackage[a4paper,margin=1.9cm]{geometry}
\usepackage[T1]{fontenc}
\usepackage{textcomp}
\usepackage[spanish]{babel}
\usepackage{amsmath,amssymb}
\usepackage{lmodern}
\renewcommand{\familydefault}{\sfdefault}
\usepackage{enumitem}

\newcommand{\vect}[2]{\begin{pmatrix}#1\\#2\end{pmatrix}}
\usepackage{tikz}
\usetikzlibrary{calc,arrows.meta,patterns,decorations.pathmorphing}
\usepackage[table]{xcolor}
\usepackage{multicol}
\usepackage{parskip}

\definecolor{unalblue}{RGB}{31,73,124}

\newlist{opciones}{enumerate}{1}
\setlist[opciones]{label=(\alph*),itemsep=6pt,topsep=6pt,leftmargin=1.8em,font=\normalfont}
\setlist[enumerate,1]{itemsep=1.4em,topsep=1em}

\newcommand{\examheader}{%
  \noindent
  \begin{minipage}[t]{0.6\linewidth}
    {\small\bfseries Geometría Vectorial y Analítica --- Parcial 1}\\
    {\small Semestre 2015--02}
  \end{minipage}%
  \hfill
  \begin{minipage}[t]{0.35\linewidth}
    \raggedleft\small Escuela de Matemáticas. Universidad Nacional de Colombia, Sede Medellín.
  \end{minipage}\\[2pt]
  \rule{\linewidth}{0.6pt}
}
\newcommand{\examfooter}{%
  \vfill
  \rule{\linewidth}{0.4pt}\\[1pt]
  {\scriptsize\textit{Transcripción digital --- MathUNAL}\hfill\textcolor{unalblue}{\textbf{\thepage}}}
}
\pagestyle{empty}

\newcommand{\nota}[1]{{\small\itshape\color{unalblue!70!black} #1}}

\begin{document}
\examheader

\begin{center}
{\Large Primer Examen Parcial de Geometría}\\[4pt]
{\small 21 de septiembre del 2015}
\end{center}

\vspace{10pt}
\textbf{PUNTAJE. SOLO PARA USO OFICIAL}

\begin{center}
\begin{tabular}{|c|c|c|c|c|c|}
\hline
1 (/30) & 2 (/30) & 3-6 (/24) & 7 (/16) & Total & Nota \\
\hline
 & & & & & \\
\hline
\end{tabular}
\end{center}

\vspace{10pt}
\textbf{INSTRUCCIONES}

Antes de empezar verifique que su examen esté completo y que sus datos sean correctos. Por favor verifique que su celular esté apagado. No está permitido el uso de calculadora ni de otros aparatos electrónicos. Solucione las preguntas en los espacios en blanco asignados a cada una de ellas. No está permitido sacar hojas en blanco ni ningún tipo de apuntes durante el examen. \textbf{Duración: 1 hora y 50 minutos.}

\vspace{6pt}
\nota{En cada uno de los literales de los puntos 1 y 2, debe presentar detalladamente el procedimiento para llegar a la solución. NOTA: Respuestas sin procedimiento no se tendrán en cuenta.}

\begin{enumerate}
\item Considere la recta $\mathcal{L}$ con ecuación general $6x-2y+1=0$.

\begin{enumerate}[label=(\alph*)]
\item{}[10pt] Encuentre una ecuación vectorial paramétrica para la recta $\mathcal{L}'$ que es paralela a $\mathcal{L}$ y pasa por el punto $P=\vect{1}{0}$.
\item{}[10pt] Encuentre la proyección ortogonal de $P$ sobre $\mathcal{L}$.
\item{}[10pt] Halle la distancia entre $\mathcal{L}$ y $\mathcal{L}'$.
\end{enumerate}

\item Considere la recta $\mathcal{L}$ que pasa por los puntos $P=\vect{-1}{0}$ y $Q=\vect{-2}{\sqrt{3}}$.

\begin{enumerate}[label=(\alph*)]
\item{}[10pt] Halle una ecuación en la forma general para la recta $\mathcal{L}$.
\item{}[10pt] Encuentre el ángulo de inclinación y la pendiente de la recta $\mathcal{L}$.
\item{}[10pt] Halle una ecuación en la forma general para la mediatriz del segmento $\overline{PQ}$. Recuerde que la mediatriz de un segmento es la línea recta perpendicular al segmento y que pasa por su punto medio.
\end{enumerate}
\end{enumerate}

\vspace{6pt}
\nota{En las preguntas 3 a 6 complete los espacios en blanco o elija la (las) opción (opciones) correcta (correctas) según el caso. Marque con una X la letra de su elección. NOTA: En esta sección se califica sólo la respuesta y no se tiene en cuenta el procedimiento.}

\begin{enumerate}
\setcounter{enumi}{2}
\item{}[4pt] Considere los vectores $X=\vect{y}{1}$ y $Y=\vect{9}{y}$. ¿Para cuáles valores de $y$ son linealmente dependientes los vectores $X,Y$?

Respuesta \underline{\hspace{4cm}}

\item{}[4pt] Suponga que $P=\dfrac{3}{8}A+\dfrac{5}{8}B$. ¿Cuál de los seis puntos del segmento de la figura es igual a $P$?

\begin{center}
\begin{tikzpicture}[scale=1]
\draw[thick] (0,0) -- (8,0);
\foreach \x/\lbl/\pos in {0/B/below,1.6/X/below,3.2/Y/below,4.8/Z/below,6.4/W/below,8/A/below}{
  \draw[thick] (\x,0.08) -- (\x,-0.08);
  \node[\pos] at (\x,-0.15) {$\lbl$};
}
\end{tikzpicture}
\end{center}

Respuesta $P=$ \underline{\hspace{3cm}}

\item{}[4pt] Considere el segmento dirigido de la figura. Se sabe además que $\|A-B\|\neq 0$ y que $\|A-B\|\neq 1$. ¿Cuál o cuáles de las siguientes son parametrizaciones de dicho segmento?

\begin{center}
\begin{tikzpicture}[scale=1]
\draw[-{Latex}] (0,0) node[below]{$A$} -- (3,2) node[above]{$B$};
\end{tikzpicture}
\end{center}

\begin{opciones}
\item $(1-t)B+tA,\quad 0\le t\le 1.$
\item $(1-t)A+tB,\quad 0\le t\le 1.$
\item $B+\dfrac{t}{\|A-B\|}A,\quad 0\le t\le \|A-B\|.$
\item $B+\dfrac{t}{\|A-B\|}A,\quad 0\le t\le 1.$
\end{opciones}

\item{}[12pt] En el rectángulo $OBCE$ se definen las rectas $\mathcal{L}$ y $\mathcal{L}'$ como se muestra en la figura. Se sabe que $APBQ$ y $AQOR$ son cuadrados. Establezca si las siguientes afirmaciones son verdaderas o falsas.

\begin{center}
\begin{tikzpicture}[scale=0.9]
\coordinate (O) at (0,0);
\coordinate (B) at (6,0);
\coordinate (C) at (6,4);
\coordinate (E) at (0,4);
\coordinate (A) at (2.4,1.6);
\coordinate (Q) at (3.6,3.6);
\coordinate (P) at (5.6,3.4);
\coordinate (R) at (0.4,3.4);
\coordinate (S) at (1.6,4.6);
\draw[thick] (O) -- (B) -- (C) -- (E) -- cycle;
\draw (A) -- (Q) -- (P) -- (B) -- cycle;
\draw (A) -- (Q) -- (R) -- (O) -- cycle;
\draw[dashed] (O) -- (P) node[midway,above]{};
\draw[dashed] (R) -- (B);
\node[below left] at (O) {$O$};
\node[below right] at (B) {$B$};
\node[above right] at (C) {$C$};
\node[above left] at (E) {$E$};
\node[below] at (A) {$A$};
\node[above] at (Q) {$Q$};
\node[right] at (P) {$P$};
\node[left] at (R) {$R$};
\node[above] at (S) {$S$};
\node[left,font=\small] at ($(O)!0.5!(P)$) {$\mathcal{L}$};
\node[right,font=\small] at ($(R)!0.5!(B)$) {$\mathcal{L}'$};
\end{tikzpicture}
\end{center}

\begin{enumerate}[label=(\alph*)]
\item $\operatorname{dist}(B-A,\mathcal{L}')=\|B-C\|$. \hspace{1cm} (V)\hspace{4pt}(F)
\item $\mathrm{P}_{\mathcal{L}}(R-B)=R$. \hspace{1cm} (V)\hspace{4pt}(F)
\item $(Q-A)\cdot(B-A)=0$. \hspace{1cm} (V)\hspace{4pt}(F)
\end{enumerate}

\item Considere el triángulo $\triangle ABC$ de la figura. Sean $R,P,Q$ los puntos medios de los segmentos $\overline{AB}$, $\overline{BC}$ y $\overline{CA}$ respectivamente.

\begin{center}
\begin{tikzpicture}[scale=0.9]
\coordinate (A) at (5,0);
\coordinate (B) at (0,1.6);
\coordinate (C) at (4,3.6);
\coordinate (R) at ($(A)!0.5!(B)$);
\coordinate (P) at ($(B)!0.5!(C)$);
\coordinate (Q) at ($(C)!0.5!(A)$);
\draw (A) -- (B) -- (C) -- cycle;
\draw[dashed] (R) -- (P) -- (Q) -- cycle;
\node[right] at (A) {$A$};
\node[left] at (B) {$B$};
\node[above] at (C) {$C$};
\node[below] at (R) {$R$};
\node[left] at (P) {$P$};
\node[above right] at (Q) {$Q$};
\end{tikzpicture}
\end{center}

\begin{enumerate}[label=(\alph*)]
\item{}[6pt] Exprese los vectores $P-R$ y $Q-R$ en términos de $A,B$ y $C$.
\item{}[10pt] Suponga que $\triangle ABC$ es un triángulo rectángulo con ángulo recto en el vértice $C$. Pruebe que el triángulo $\triangle PQR$ es un triángulo rectángulo con ángulo recto en el vértice $R$.
\end{enumerate}

\end{enumerate}

\examfooter
\end{document}
